\section{Introduction}

\PARstart{C}{ompilers} apply a sequence of optimization passes, not a single optimization, and the passes depend on one another. One pass can
create the pattern that a later pass needs, remove information that another pass relies on, or change a cost that a later
decision reads. This is the phase-ordering problem. Its size on real code is seldom measured directly. Most published experiments
report which sequence performs best, not how much each pair of passes changes the result.

This paper measures those interactions in the new pass manager of LLVM \cite{lattner2004llvm}. We ran \NumSweeps{} pipeline variants over \NumKernels{}
kernels from four benchmark suites. The main measure is the instruction count of the optimized IR, which is exact and independent
of timing and of the backend. We also timed every variant against the canonical build of its kernel in one measurement window, and
we checked every variant's optimized IR against its canonical IR. The results are stored in a telemetry table with \NumColumns{}
fields per sweep. We ask four questions.

\begin{enumerate}
  \item How much does the order of two passes change the code, and which pairs interact most? (Section~\ref{sec:results-instcount}.)
  \item How much of that change is visible in the execution time, given the noise of the timing measurements? (Section~\ref{sec:results-timing}.)
  \item How many pass executions in a standard pipeline leave the IR unchanged, and what time do those executions consume? (Section~\ref{sec:results-rpb}.)
  \item How far can a bounded, machine-checked comparison of optimized code go towards showing that each transformation preserves
        meaning? (Section~\ref{sec:results-refinement}.)
\end{enumerate}

Our contributions are the following.
\begin{itemize}
  \item An empirical campaign of \NumSweeps{} sweeps. It covers the single passes, every ordered pair of \NumThreePasses{} passes,
        a stepwise ablation of the \texttt{-O3} pipeline, an interference study around \texttt{loop-vectorize}, and sequences chosen
        from the instruction-count tensor.
  \item An instruction-count interaction, normalized by the canonical count, and an identity for self-interactions that holds
        exactly in all \InstSelfIdentityCases{} kernel-pass cases where it applies.
  \item IR-level checks of every pass on two canonical forms, which show that the canonical form changes the effect of
        \texttt{loop-unroll}.
  \item A timing measure with its noise floor, and the finding that timing resolves few pass interactions, all of them below the noise floor.
  \item A block-local refinement checker for integer LLVM IR. We state the soundness argument for its UNSAT answers, we check its
        encoding exhaustively against an independent reference semantics for small widths, and we report where its answers stop
        being evidence of correctness.
\end{itemize}

